\section{My team}

\subsection*{2a. Structure --- discussed, drifted, re-opened}

At the kickoff around wk13 we did have an explicit structure conversation, and I want to give the team credit for that because it is a thing I have failed to do in past projects. We split on module ownership: Robin took flight and state-machine execution, I took perception and simulation, Edward took the Cube and hardware interface, Demetro took client communication and presentations, and [PLACEHOLDER: PM name] nominally coordinated. We did not name a single technical lead, which was a deliberate choice --- the brief (R08) gives teams discretion over autonomy levels and interface design, and we wanted that discretion to be distributed.

The structure drifted twice. The first drift was silent: around wk16, as hardware kept slipping, the centre of gravity of the project quietly migrated from the bench to my laptop, because that was the only place where forward progress was possible. Nobody discussed it, nobody declared it, it just happened. The second drift was louder. Around wk19, when it became clear Robin needed a cleanly-scoped flight script that did not try to also do vision, we sat down and explicitly renegotiated the interface: his script would handle the flight loop and state transitions, and call into my \texttt{passive\_watch} as an HTTP service for detection results. I wrote this down in \texttt{docs/DESIGN\_DECISIONS.md} (lines 107--109) specifically because I had watched the first drift happen without anyone noticing, and I did not want the second one to.

In hindsight, that second explicit renegotiation is one of the things I would do earlier next time. The first silent drift cost us about two weeks of parallel work where Robin and I were building overlapping state machines without realising it, and I only noticed the overlap when I read his code in preparation for the complaint letter. What looked like a technical duplication was actually a communication failure, because we had never agreed out loud on who owned what when the hardware disappeared. I now see that unexplicit structure is a luxury that only survives in teams with a lot of slack. Under hardware pressure, structure must be restated every time the environment shifts. I did not know this before this project. I know it now.

\subsection*{2b. What was most impactful --- and one episode of conflict}

The single most impactful element of team working, honestly, was the complaint letter process. Not because of what we said to Steve --- the letter's eventual impact on our grade remains to be seen --- but because of what drafting it did \emph{inside} the team. I had been sitting on a private frustration for weeks. Robin had been sitting on his. Demetro had been sitting on his. We had been absorbing the hardware failures individually, politely, and in parallel, because none of us wanted to be the one who complained. Getting everyone on a call to review the first draft was the first time the team actually talked about how bad the situation was. That conversation shifted something in our working rhythm --- after it, we stopped pretending that hardware would arrive next week and started planning as if it would not.

There was also one specific conflict I want to name, because pretending it did not happen would be dishonest. Around wk17, Robin and I disagreed about the architecture of the flight state machine. I had pushed for a stricter finite-state machine with explicit transitions and timeout guards; Robin preferred a simpler polling loop with fewer states. I initially framed my position as a technical argument --- I genuinely believed the FSM was safer because of the geofence and NFZ edge cases --- but in hindsight I can see at least half of my insistence was because the FSM was my code and I was emotionally invested in it. Robin, to his credit, did not escalate. He pushed back clearly, once, in a Teams message, saying roughly that my FSM was over-engineered for a codebase that two out of five teammates could read. I took 24 hours to calm down, reread his flight script, and realised he was partly right --- my FSM had 32 transitions, and nobody except me was going to maintain it. We compromised: I kept the FSM for the full autonomous path in \texttt{main.py} but stripped it down to a lightweight polling structure for the passive-watch integration that Robin's script calls. That compromise is the reason commit \texttt{a5b1ab7} exists: the \texttt{--robin} flag is not a feature, it is an interface concession.

The lesson was not that I should defer on technical disagreements --- I still think the FSM is the right structure for the full mission --- but that I should separate ``is this the right design'' from ``is this the right design \emph{for this team}.'' A design that only one person can maintain is not the right design for a five-person team with one flight day left. I was right about the FSM and wrong about the context. I had not previously noticed that these two kinds of rightness could come apart in me.

\subsection*{2c. What I would change in future}

Three things, in order of importance.

First, I would \textbf{institute a weekly decision log} from day one. Not a meeting minute --- a one-line record of every call that bound more than one person's work, with a date, an owner, and a reason. We lost at least a week in wk16--wk17 to drift that would have been caught by any written record. I already started doing this informally in \texttt{NICE\_TO\_HAVE.md} and \texttt{brain-dump.md} for my own work, and the effect on my ability to onboard a new session was measurable. Scaling it to the team level is the single highest-leverage process change I can name.

Second, I would \textbf{front-load interface contracts before parallel work begins}. The Robin-vs-me state-machine overlap in wk16 was exactly the failure mode the M16 guide warns about: high-agency individuals default to parallel solo work under stress, which hides disagreements until they become expensive. A two-hour whiteboard session in wk14 to agree on the HTTP interface between the flight script and the vision script would have saved the integration week in wk20 entirely.

Third, I would \textbf{stop assuming that documentation is a substitute for conversation}. I wrote a lot of markdown files for Robin, Demetro, and a hypothetical future colleague. Most of them were read partially or not at all. The two pieces of writing that actually changed teammates' behaviour were (a) the Teams message I sent to the group after the wk18 cancellation proposing the simulation pivot, which was 400 words, and (b) the complaint letter, which was 2000 words but included a live group-call discussion before sending. Both worked because they were short, timely, and accompanied by a conversation. The 3000-word VISION\_PIPELINE\_FOR\_COLLEAGUE doc, by contrast, sat in the repo and was never used, because I never walked anyone through it. The lesson is not that documentation is useless --- it is that documentation without a conversation is a message in a bottle, and the asynchronous-purist in me needs to stop using it that way.
